% page 535 of the facsimile (image p-534 of the scan)
% block 188.0 x 246.3 mm ; left margin 14.0 mm
\pgc{-7.620mm}{0.000mm}{
\l{\cel{59}527.}
\l{}
\l{\cel{6}signatar.\cel{2}(\sou{trans.}) Skribar sua nomo sub la texto di letro, di}
\l{\cel{9}akto, e c., kom atesto pri to ke lu ipsa skribis ta texto,}
\l{\cel{9}o ke lu aprobas olca. - DEFIS.}
\l{}
\l{\cel{6}signalar.\cel{2}(trans., ad)\cel{3}Indikar per signo ke ulo esas agenda}
\l{\cel{9}ye instanto determinita. DEFIS.}
\l{}
\l{\cel{6}signifikar. (\sou{trans.}) (\sou{aludante senco}). Enumerar la objekti a qui}
\l{\cel{9}la senco di vorto o vorti esas aplikebla. (Ex.: La vorti \textquotedbl{}la}
\l{\cel{9}maxim granda de la nombri\textquotedbl{} havas senco komprenebla e mem klara,}
\l{\cel{9}ma korespondas a nula objekto ; lu signifikas nulo. \textquotedbl{}Blanka}
\l{\cel{9}Monto\textquotedbl{} havas senco aplikebla ad omna blanka monti ; ma lua}
\l{\cel{9}signifiko esas restriktita ad un monto.(Cf.\cel{1}\sou{Progres}o, marto}
\l{\cel{9}l9l3, pag. 80 e 8l, artiklo da Louis Couturat).- EFIS.}
\l{}
\l{\cel{6}\sou{sika}. Qua ne plus kontenas en sua korpo o sur sua surfaco parti-}
\l{\cel{9}kuli aqua. - EFIRS.}
\l{}
\l{\cel{6}\sou{siklo}. Mancheto ligna kun ferajo mi-cirkla, tranchiva (kun o sen}
\l{\cel{9}denti), quan onu uzas por mesonar la stipi di frumento, di se-}
\l{\cel{9}kolo, e c. - DE.}
\l{}
\l{\cel{6}\sou{sikomoro}. (\sou{bot.)}\cel{2}Varietato di figiero di qua la folii pasable}
\l{\cel{9}similesas olti di morusiero. - L.\cel{1}\sou{ficus sycomorus}. DEFIS.}
\l{}
\l{\cel{6}\sou{silo}.\cel{2}Exkavajo, facita en sulo, aden qua onu depozas la grani}
\l{\cel{9}(di frumento, e c.) por konservar li. - DEFIS.}
\l{}
\l{\cel{6}\sou{silabo}. (\sou{fonetiko}). Son-uno artikulata, produktata per un voc-}
\l{\cel{9}emiso, e qua konsistas ye o konsonanti e vokali, o vokalo}
\l{\cel{9}sola, o diftongo. - DEFIRS.}
\l{}
\l{\cel{6}\sou{sileno.(bot.)}\cel{2}Herbo un-yara, du-yara o perena, ed, en ula kazi,mi-}
\l{\cel{9}arboreto di qua la folii interopozesas, e kun flori (multa-}
\l{\cel{9}kaze) grupigita aden +kimi.\cel{2}L.\cel{1}\sou{silenus}.}
\l{}
\l{\cel{6}\sou{silencar.}\cel{2}(\sou{netrans.})\cel{2}Nule bruisar. - EFIS.}
\l{}
\l{\cel{6}\sou{silepso}. (\sou{retor}.)Figuro qua konsistas ye uzar vorto (che la lingui}
\l{\cel{9}en qui existas genro qua ne relatas unike la sexuozi) sen egar-}
\l{\cel{9}dar la genro, nek la nombro gramatikala, di altra vorto kun qua}
\l{\cel{9}lu relatas - omisante la konkordo gramatikala pro egardar unike l}
\l{\cel{9}ideo ipsa. - DEFIRS.}
\l{}
\l{\cel{6}\sou{silexo}. (\sou{mineral.})\cel{2}Stono tre harda, qua konsistas esence ye}
\l{\cel{9}silico, e qua produktas cintili se lu shokesas da stalo-peco.}
\l{\cel{9}- EFIS.}
\l{}
\l{\cel{6}\sou{silfo}. (\sou{mitol}.)Ento fantastika, feo di aero. - DEFIRS.}
\l{}
\l{\cel{6}\sou{siliko}. (\sou{kemio}). Korpo simpla, metala, qua, kombinite kun oxo -}
\l{\cel{9}formacas silico. -\cel{1}\sou{Simbolo kemiala}\cel{1}:Si. - DEFIS.}
\l{}
\l{\cel{6}silico. (\sou{kemio}). Oxido di siliko, qua konstitucas silexo.\cel{1}\sou{Simbolo}}
\l{\cel{9}kemiala :\cel{2}Si O2 . EFIS.}
}
